% problem-000361 7 晶体结构与晶胞
\section*{7. 晶体结构与晶胞}
\textit{下列为分问。}

\subsection*{7-1}
Si4+、Al3+ 和 $\mathrm { O ^ { 2 \cdot } }$ −的离⼦半径分别为 41 pm、50 pm 和 140 pm，通过计算说明在⽔合铝硅酸盐晶体中Si4+和Al3+各占据由氧构成的何种类型的多⾯体空隙。

\subsection*{7-2}
上述化学式中的�等于多少？说明理由。若M为2价离⼦，写出�与�的关系式。

\subsection*{7-3}
X光衍射测得Si—O键键⻓为160pm。此数据说明什么？如何理解？

\subsection*{7-4}
说明以下事实的原因：

① 硅铝⽐（有时⽤ $\mathrm { S i O _ { 2 } / A l _ { 2 } O _ { 3 } }$ 表⽰）越⾼，分⼦筛越稳定；② 最⼩硅铝⽐不⼩于 $1 _ { \circ }$

\subsection*{7-5}
⼈⼯合成的A型分⼦筛钠盐，属于⽴⽅晶系，正当晶胞参数 $a = 2 4 6 4 \mathrm { p m }$ ，晶胞组成为

$\mathrm { N a _ { 9 6 } [ A l _ { 9 6 } S i _ { 9 6 } O _ { 3 8 4 } ] { \cdot } { \cdot } { \cdot } { \cdot } { \mathrm { ~ H _ { 2 } O _ { \circ } ~ } } }$ 将811.5克该分⼦筛在 $1 . 0 1 3 2 5 \times 1 0 ^ { 5 } \mathrm { P a } , ~ 7 0 0 ~ ^ { \circ } \mathrm { C }$ 加热6⼩时将结晶⽔完全除去，得到798.6升⽔蒸⽓（视为理想⽓体）。计算该分⼦筛的密度 $\dot { \boldsymbol { D } } _ { \mathrm { c } }$ 。
